\documentclass[11pt]{article}
\usepackage[margin=1in]{geometry}
\usepackage{amsmath,amssymb,amsthm,mathtools,booktabs,hyperref}
\newtheorem{definition}{Definition}
\newtheorem{proposition}{Proposition}
\newtheorem{conjecture}{Conjecture}
\title{Scarcity-Amplifying Institutional Feedback, Sampled Governance,\\and Safety-Relevant Nonlinear Transitions}
\author{Author(s) redacted for review}
\date{}
\begin{document}\maketitle

\begin{abstract}
This paper studies a restricted family of institutional-feedback stress-test models in which perceived scarcity can amplify extraction rather than trigger protective restraint. It does not claim that institutions generally behave this way, that continuous delay represents actual governance by default, or that a model-specific bifurcation threshold is a universal policy threshold. The paper compares scarcity-amplifying, protective, and inertia/capture response structures; embeds policy in a sampled hybrid review process; distinguishes stock-culling from recruitment-suppression liquidation; and defines a nonlinear transition as sustainability-relevant only when a reachable invariant set, basin, or uncertainty tube intersects a declared unsafe set. Existing scalar delay equations, gated variants, stage structures, and four-state support-pool extensions are retained as model-specific sensitivity families subject to independent numerical reproduction and empirical identification.
\end{abstract}

\section{Scope and status}
This paper preserves the original nonlinear feedback programme but changes its interpretation. A Hopf crossing, periodic-orbit fold, or large cycle is a property of a specified dynamical model. It becomes a sustainability result only after the model has a typed physical state, an identified observation/policy rule, a declared safe set, and evidence that the relevant invariant set or reachable tube approaches or crosses that safe set.

The paper uses three policy hypotheses:
\[
\text{H1: scarcity-amplifying extraction},\qquad
\text{H2: protective restraint/restoration},\qquad
\text{H3: inertia/capture/state-dependent response}.
\]
No result for H1 is generalized to H2 or H3.

\section{Physical liquidation channels}
Let $N$ denote an immediately harvestable standing stock, $J$ a recruitment/immature class when relevant, and $E$ extraction effort. The phrase ``liquidation'' has at least three distinct meanings:
\begin{enumerate}
\item standing-stock culling: present extraction removes reproductive stock directly;
\item recruitment suppression: present use prevents future recruits without immediate adult removal;
\item weak viability coupling: use has limited or indirect effect on reproduction.
\end{enumerate}
A diagnostic label such as ``unsustainable portion'' never determines physical destination. Material routing is determined by the typed physical module, not by whether a flow exceeds a service or regeneration benchmark.

\section{State space, units, and solution status}
For an RFDE variant, the state is a history in a declared space such as $C([-\tau_{\max},0],\mathbb R_+^3)$, together with any separately typed discrete institutional mode. Initial histories, parameter domains, local Lipschitz/Carath\'eodory conditions, continuation assumptions, and the solution concept are part of the model. The three coordinates have distinct units: $N$ is stock, $Z$ has the units of the declared filtered deficit, and $E$ is an effort/deployment intensity; every coefficient must be dimensioned accordingly. Positivity and an effort bound are theorems only after the boundary vector field, gates, and resets are checked.

\section{Continuous-time stress-test family}
A minimal stress-test family is
\begin{align}
\dot N&=R(N,A,\ldots)-qE N,\label{eq:N}\\
\dot Z&=\tau_m^{-1}[D(N,E,A,\ldots)-Z],\label{eq:Z}\\
\dot E&=\mathcal G_\gamma(E,Z(t-\tau_I),h),\label{eq:E}
\end{align}
where $Z$ is filtered perceived scarcity/deficit and $\gamma\in\{\mathrm{H1,H2,H3}\}$ indexes the policy mechanism. The equations are a phenomenological stress-test family unless derived from a domain module. $A$ may represent a support pool only when its material dynamics are included; a frozen $A$ is a formal limiting case, not a default physical assertion.

For H1, $\partial\mathcal G_{\rm H1}/\partial Z>0$ over the relevant range; for protective feedback, the effect on extraction is reversed or includes restoration controls. The model must state whether $E$ is an endogenous industry response, a legal quota-utilization state, or an actual institutional control; these interpretations are not interchangeable.

\paragraph{Non-negativity and effort bounds.} Every variant must verify present-donor stock non-negativity, effort admissibility, and behavior at all boundaries. Gating, damping, or saturation is not accepted merely because it preserves a desired bifurcation; it must have a constitutive interpretation and an independently stated parameter role.

\section{Sampled governance as the default institutional representation}
At review times $t_k$, institutions observe data and update prescriptions:
\[
Y_k=\mathcal O(X_{[t_{k-1},t_k]})+\varepsilon_k,
\qquad
a_{k+1}=\Pi_\gamma(Y_k,h_k,\rho_k).
\]
The implemented control satisfies $U_{k+1}\in\mathcal E(B_k,h_k,a_{k+1})$ and is held by the registered zero-order-hold rule until the next review. Any interpolation or dispatch queue defines a separately identified controller. This hybrid system is the primary governance representation.

A continuous-delay approximation to \eqref{eq:E} is admissible only after specifying the review interval, hold rule, implementation lag, and approximation error. A sampled map can exhibit flip, Neimark--Sacker, border-collision, or other discrete-time behavior not captured by continuous-delay Hopf terminology.

\section{Alternative biological mechanisms}
The stress-test programme retains alternatives that can mimic or replace an institutional cycle:
\begin{itemize}
\item stage/maturation delays and cohort resonance;
\item recruitment suppression versus stock culling;
\item support-pool/nutrient limitation;
\item pollution/waste feedback;
\item predator--prey or climatic forcing;
\item direct abiotic mining and slow-store exhaustion.
\end{itemize}
A candidate institutional mechanism must be compared with these alternatives rather than identified from periodicity alone.

\section{Safety-relevant nonlinear transitions}
Let $\mathcal C_X$ be the state-safe set. A local bifurcation in a reduced model is not automatically a sustainability transition. Distinguish:
\begin{enumerate}
\item existence of a local bifurcation in a specified reduced model;
\item persistence of the organizing structure under physically admissible coupling;
\item safety relevance, measured by intersection of an attractor, basin, or uncertainty tube with $\mathcal C_X^c$.
\end{enumerate}
Useful quantities include
\[
\operatorname{dist}(\mathcal A_\mu,\mathcal C_X^c),
\qquad
\operatorname{dist}(\mathcal R_{[0,T]},\mathcal C_X^c),
\]
where $\mathcal A_\mu$ is an invariant set and $\mathcal R_{[0,T]}$ a reachable tube. Basin membership, not only attractor location, matters in bistable systems.

\begin{conjecture}[Structured persistence]
A nonlinear transition verified in a reduced hybrid/RFDE model may persist under small, physically admissible coupling only after well-posedness, fast difference-operator contractivity, spectral separation, center-manifold reduction, transverse Poincar\'e-map conditions, and preservation of positivity/safety are established. At a periodic-orbit fold, transverse normal hyperbolicity does not imply normal hyperbolicity of the fold orbit itself.
\end{conjecture}

\section{Variant registry and numerical standards}
The archived programme contains ungated, gated, hybrid-effort, four-state support-pool, two-channel liquidation, stage-structured, sampled-review, thermodynamic-tether, and unified-core variants. Every retained result must be tabulated by:
\begin{itemize}
\item governing equations and state/control interpretation;
\item parameter units and identifiability status;
\item equilibrium and boundary conditions;
\item numerical method, mesh/time-step, event handling, and horizon;
\item branch identity and uncertainty;
\item status: analytical, independently reproduced numerical, conjectural, or superseded.
\end{itemize}
A persistence boundary, a fold, an SNPO-confirmed fold, and an initial-condition basin observation are different claims. No fold is called SNPO-confirmed without branch-specific multiplier/transversality evidence. Safety-relevant numerics additionally require convergence/error evidence for delay discretization, event handling, orbit and Floquet calculations, positivity constraints, and distance of the attractor/reachable tube to the unsafe set. Where feasible, use independent solvers and validated/interval enclosures rather than a single continuation implementation.

\section{Empirical identification and case screening}
The empirical programme begins with a preregistered candidate universe and screening protocol. A candidate must supply: individual-resource data; dated observation, assessment, authorization, and implementation events; an identified policy sign; plausible review/implementation dynamics; major alternative mechanisms; and enough record length relative to the predicted phenomenon.

Failure to identify a clean case is reported as a limitation of the sample/protocol, not a general empirical null, unless the candidate universe and exclusion logic justify that inference. Cross-sectional annual-management nulls are descriptive only unless a causal identification design supports stronger claims.

\section{Policy scope}
This paper does not advise slowing protective governance. It asks whether, under an identified scarcity-amplifying response law, faster mobilization can intensify positive feedback. Protective quotas, restoration, monitoring, authority, and implementation may have different signs and timing effects. Policy recommendations require hybrid MSE under stated uncertainty, authority, compliance, distributional, and safety constraints.

\section{Conclusion}
The original delay-bifurcation work is preserved as a valuable but narrow stress-test programme. Its scientific contribution is conditional: it can reveal how specific positive scarcity-to-extraction feedbacks, review rules, biological lags, and nonlinearities may generate unsafe dynamics. It is not a universal engine of sustainability or a substitute for typed physical, observational, and institutional modules.

\end{document}
